\section{Discussion}
\label{sec:discussion}

\heading{Non-determinism in skill compilation.}
Unlike traditional program compilation, skill compilation takes natural language as input,
which inherently introduces some non-determinism into the compilation process~\cite{wang2023selfconsistency}.
However, the LLM executing skills differs fundamentally from traditional hardware such as CPUs~\cite{hennessy2019architecture},
and possesses an inherent tolerance for input variability that enables skills to execute correctly.
Moreover, \sys's rigorous compilation optimization passes and rollback mechanisms ensure
that compiled skills consistently deliver stable performance improvements across most downstream tasks.


\heading{Capability Coverage.}
The current catalog has 26 primitive capabilities across four categories,
enough to cover the 15,063 skills we analyzed. 
As the skill library grows, we can adopt the iterative approach described in \S\ref{sec:compiler:cap} to expand the primitive capabilities.

\heading{Compilation cost.}
AOT compilation invokes an LLM, which incurs token costs. 
However, since skills are executed repeatedly at runtime,
the compilation overhead is amortized across multiple invocations~\cite{gray1981transaction}.
Furthermore, compiled skills can be shared across users and applications, further reducing the per-skill compilation cost.
